% BHU PYQ -- Transcribed & Corrected
% Paper: CHEMJ42: Organic Chemistry
% Exam: B.Sc. (Hons) Semester IV Examination 2025-26 (NEP)
% Ref Code: Conf/Even/N/2025-26/06
% Category: Major
% Department: Chemistry

\documentclass[11pt,a4paper]{article}
\usepackage[a4paper,top=2.5cm,bottom=2.5cm,left=2.8cm,right=2.8cm,headheight=14pt]{geometry}
\usepackage{amsmath,amssymb}
\usepackage[T1]{fontenc}
\usepackage[utf8]{inputenc}
\usepackage{lmodern,microtype,fancyhdr,enumitem,hyperref,lastpage,parskip}

\hypersetup{
    colorlinks=true,
    urlcolor=black,
    linkcolor=black,
    pdftitle={CHEMJ42: Organic Chemistry -- BHU Sem IV 2025-26 (NEP)}
}

\pagestyle{fancy}
\fancyhf{}
\renewcommand{\headrulewidth}{0.4pt}
\renewcommand{\footrulewidth}{0.4pt}
\fancyhead[L]{\small \textit{Banaras Hindu University}}
\fancyhead[R]{\small \textit{CHEMJ42: Organic Chemistry (2025-26)}}
\fancyfoot[C]{\small Page \thepage\ of \pageref{LastPage}}

\newcommand{\pts}[1]{\hfill \textit{[#1]}}

\begin{document}

\begin{center}
  {\large\bfseries BANARAS HINDU UNIVERSITY}\\[2pt]
  {\normalsize B.Sc. (Hons) Semester IV Examination 2025-26 (NEP)}\\[6pt]
  \rule{0.8\linewidth}{0.5pt}\\[6pt]
  {\large\bfseries CHEMISTRY}\\[4pt]
  {\large\bfseries Paper Code: CHEMJ42 : Organic Chemistry}\\[2pt]
  {\small Ref. No.: Conf/Even/N/2025-26/06}\\[4pt]
  {\normalsize Time: 2$\frac{1}{2}$ Hours\hfill Maximum Marks: 70}
\end{center}

\rule{\linewidth}{0.4pt}

\smallskip

\noindent \textit{Instructions: Answer five questions including Question 1 which is compulsory. The figures in the right-hand margin indicate marks.}

\bigskip

\noindent \textbf{Question 1.} Answer the following: \pts{7 $\times$ 2 = 14}
\begin{enumerate}[label=(\alph*)]
  \item Write down the cis- and trans-isomers of indigotin and explain their stability.
  \item Give two suitable methods for the protection of amino ($-\text{NH}_2$) group in organic synthesis.
  \item Write resonating structures of thiophene using d-orbitals.
  \item Explain anomeric effect with suitable example.
  \item Define glycosylation reaction giving suitable example.
  \item Describe any one method for the synthesis of cyclopropane.
  \item Complete the following reactions:
  \begin{enumerate}[label=(\roman*)]
    \item $\text{Furan} + \text{NH}_3 \xrightarrow[\text{heat}]{\text{Al}_2\text{O}_3, \text{H}_2\text{O}} ?$
    \item $\text{Indole} + \text{HNO}_3, \text{Ac}_2\text{O} \rightarrow ?$
  \end{enumerate}
\end{enumerate}

\bigskip

\noindent \textbf{Question 2.}
\begin{enumerate}[label=(\alph*)]
  \item Give the chemical equations and names of the product formed, when nitrobenzene reacts with--- (i) $\text{Sn/HCl}$ (ii) $\text{Zn/NH}_4\text{Cl}$ (iii) Methanolic $\text{NaOH}$ and zinc dust (iv) $\text{Zn/CH}_3\text{OH/NaOH}$. \pts{4 $\times$ 1 = 4}
  \item Explain the term epimerization using a suitable example. Discuss the interconversion of D-glucose and D-mannose, providing the chemical reactions involved in the process. \pts{4}
  \item Starting from aniline, how would you prepare 1,3,5-tribromobenzene? Explain why direct bromination of benzene is not suitable for this synthesis. \pts{4}
  \item An organic compound [A] with molecular formula $\text{C}_6\text{H}_7\text{N}$ on treatment with $\text{NaNO}_2$ and $\text{HCl}$ at low temperature give [B]. Compound [B] reacts with cuprous cyanide to give [C] ($\text{C}_7\text{H}_5\text{N}$). Upon acid hydrolysis, [C] yields a carboxylic acid [D] ($\text{C}_7\text{H}_6\text{O}_2$). Identify [A] to [D] and write the reactions. \pts{2}
\end{enumerate}

\bigskip

\noindent \textbf{Question 3.}
\begin{enumerate}[label=(\alph*)]
  \item Describe the Haworth synthesis of naphthalene in detail, including all reagents and steps involved. \pts{5}
  \item Explain the synthesis of indigotin (indigo) from anthranilic acid. Also describe the vat dyeing process of indigo with suitable reaction. \pts{4}
  \item Predict the products when naphthalene reacts with--- (i) $\text{CrO}_3$ in acetic acid (ii) $\text{Na/C}_2\text{H}_5\text{OH}$ (iii) Conc. $\text{H}_2\text{SO}_4$ at $160\ ^\circ\text{C}$. \pts{3 $\times$ 1 = 3}
  \item Explain why the $\text{C}_1-\text{C}_2$ bond in naphthalene is shorter than the $\text{C}_2-\text{C}_3$ bond. \pts{2}
\end{enumerate}

\bigskip

\noindent \textbf{Question 4.}
\begin{enumerate}[label=(\alph*)]
  \item Write the product(s) in the following reactions: \pts{6 $\times$ 1 = 6}
  \begin{enumerate}[label=(\roman*)]
    \item $\text{Anthracene} + \text{Maleic anhydride} \xrightarrow{\Delta} ?$
    \item $\text{Anthrone} \xrightarrow{\Delta} ?$
    \item $\text{N-methylformanilide} + \text{POCl}_3 \xrightarrow{\text{heat}} ?$
    \item $\text{Pyridine} \xrightarrow[\text{(ii) H}_3\text{O}^+]{\text{(i) KOH}} ?$
    \item $\text{2-Furaldehyde} + \text{NaOH} \rightarrow ?$
    \item $\text{Pyridine} + \text{CH}_3\text{CO}_3\text{H} \rightarrow ?$
  \end{enumerate}
  \item Compare and discuss the aromatic character of furan, thiophene and pyrrole. Arrange them in the decreasing order of aromaticity. \pts{4}
  \item Draw all the conformations of 1,2-dimethylcyclohexane. Identify which isomer is more stable and why. \pts{4}
\end{enumerate}

\bigskip

\noindent \textbf{Question 5.}
\begin{enumerate}[label=(\alph*)]
  \item Discuss the mechanism of any three of the following: \pts{3 $\times$ 2 = 6}
  \begin{enumerate}[label=(\roman*)]
    \item Chichibabin reaction
    \item Feist-B\'enary synthesis
    \item Skraup synthesis of quinoline
    \item Fischer indole synthesis
  \end{enumerate}
  \item Describe the reaction and mechanism for synthesis of fluorescein dye. \pts{4}
  \item Write all the steps involved in the synthesis of osazone from D-glucose. Explain why glucose, fructose and mannose form the identical osazone with excess phenylhydrazine. \pts{4}
\end{enumerate}

\bigskip

\noindent \textbf{Question 6.}
\begin{enumerate}[label=(\alph*)]
  \item Discuss Baeyer's strain theory for the stability of cycloalkanes. Calculate the 'Angle strain' for cyclopropane and cyclobutane. Why does this theory fail to explain the stability of higher cycloalkanes? \pts{6}
  \item Draw the energy profile diagram for different conformations of cyclohexane. Arrange them in the decreasing order of stability and give reasons for the same. \pts{4}
  \item Describe the use of silyl ethers in protecting alcohols. How are these groups introduced and removed? \pts{4}
\end{enumerate}

\bigskip

\noindent \textbf{Question 7.}
\begin{enumerate}[label=(\alph*)]
  \item Complete the following chemical equations with appropriate structures: \pts{6 $\times$ 1 = 6}
  \begin{enumerate}[label=(\roman*)]
    \item $\text{D-Fructose} \xrightarrow{\text{NaBH}_4} ?$
    \item $\text{D-Glucose} \xrightarrow[\text{HCl}]{\text{Acetone}} ?$
    \item Methyl $\alpha$-D-glucopyranoside $\xrightarrow{\text{HIO}_4} ?$
    \item $\text{Sucrose} \xrightarrow[\text{NaOH}]{\text{(CH}_3\text{)}_2\text{SO}_4} ? \xrightarrow{\text{dil. HCl}} ?$
    \item $\text{D-Glucose} \xrightarrow[\text{H}^+]{\text{C}_2\text{H}_5\text{OH}} ?$
    \item Indigo $\xrightarrow{\text{NaOH}} ? + ?$
  \end{enumerate}
  \item How will you prepare the following from benzene diazonium chloride? \pts{3}
  \begin{enumerate}[label=(\roman*)]
    \item Biphenyl
    \item Phenylhydrazine
    \item Azobenzene
  \end{enumerate}
  \item Compare the structures of starch and cellulose. Explain the difference between $\alpha$-glycosidic and $\beta$-glycosidic linkages and how they affect the overall shape of the molecule. \pts{3}
  \item Describe the structure of sucrose. Why is it called non-reducing sugar? \pts{2}
\end{enumerate}

\bigskip

\noindent \textbf{Question 8.}
\begin{enumerate}[label=(\alph*)]
  \item Why is protection of functional groups necessary in organic synthesis? Mention the three essential criteria for an ideal protecting group. \pts{4}
  \item How will you reduce a keto group in a molecule that also contains an aldehyde group without affecting the aldehyde group? \pts{3}
  \item Explain the protection and deprotection of the following functional groups with suitable reagent: \pts{3}
  \begin{enumerate}[label=(\roman*)]
    \item Carboxyl group
    \item Terminal triple bond
    \item 1,2-diol
  \end{enumerate}
  \item Compare the basic strength of pyridine and pyrrole. \pts{2}
  \item Write the chemical equation and discuss the reaction between furan and maleic anhydride. \pts{2}
\end{enumerate}

\bigskip
\begin{center}
  \textbf{***}
\end{center}

\end{document}
